\part{Unchecked targets and remaining gates}
\label{part:open}
\section{What has not been proved}
There is no unchecked finding offered as a theorem in this part. The
following are the unresolved targets of the task. Their state is a
current gate, not a proved negative answer. A precise counterexample to
one proposed transfer is recorded separately from a proof that no other
transfer can exist.

\begin{description}
\item[O-PH2-SH: direct bounds for irreducible-inner plethysm.]
Both clauses of ordinary plethysm PH3 are now refuted: P-F19 gives
an exponential canonical parameter degree, and P-F30--P-F31 prove
actual simple poles of exponential order in the same cubic family.
The complete support classification closes the singular-boundary
and possible-cancellation gates of P-F23--P-F29. The first PH3
clause therefore cannot serve as the premise for a proof of PH2,
although the conditional implication P-F11--P-F12 remains valid.

The unresolved question is a direct proof or refutation of the
narrower polynomial positivity-index and saturation-index bounds,
with the source's exact parameter restrictions and constituent
polynomial definitions. A large quasipolynomial period does not by
itself force either index to be large. P-F32 supplies a precise
exterior-power complement identity and an infinite family with
first positive stretch $\Theta(\sqrt n)$; it gives only a polynomial
index lower bound. It does not settle SH(a), its separate density
clause SH(b), or PH2. P-F14--P-F15 already refute the broader
rank-only subgroup-restriction versions, a different formulation.

P-F34 makes a direct sufficient gate precise. Let $B_n$ be the
average constituent polynomial, determined by the principal part
of $H_n$ at one, and let $U_n$ be its coefficient $\ell^1$ norm.
For a specified primitive root, the proved Fourier amplitude satisfies
$\log_2|A_n|=n^4/24+O(n^3)$. Positivity or saturation
after a shift forces $B_n$ at that shift (or one more) to dominate
this amplitude. The additional bound
$\limsup (\log_2 U_n)/n^4<1/24$ would therefore refute polynomial
PH2 and SH(a). This bound remains unproved; the Fourier inequality
is necessary, not an equivalence to the index conjectures.
P-F42 proves an unconditional actual-family upper bound
$\limsup(\log_2 U_n)/n^4\le\log_2(5)/12$; its constant
is larger than $1/24$. Its coefficient-sensitive refinement
(P-F42z)--(P-F42ah) proves that the highest fifth of the
coefficient indices already has logarithmic leading constant
at most $\log_2(5)/60<1/24$. The sufficient norm gate is
equivalent to the same strict bound for the complementary
lower coefficients only. P-F44 identifies the exact global
primitive-Klein contribution as a signed cyclotomic trace and
gives an explicit complementary balanced-cone expression.
Its odd-$D$ contribution cancels, and even-$D$ signs can mix.
The remaining task is therefore a signed coefficient estimate,
not a bound on a single local amplitude.
The refined bound also removes the high-degree tail uniformly
on every polynomially bounded interval. Thus any polynomially
bounded strict index must make the signed lower polynomial in
(P-F42ai)--(P-F42ak) reach almost the full Fourier amplitude.
A separate route is now P-F56: it proves the exact
maximal two-part of every Hilbert-pole order and cancels all other
contributions by subtracting constituents half a period apart.
For its explicit primitive dyadic wave, every integer $x$ after a
valid strict shift must satisfy $f_n(x)\ge2W_n(x)$.
The actual coefficient has an elementary ambient-dimension bound.
The specified pair $a_*=(3g)^{-1}\pmod{2^{n-1}}$ has amplitude
at least $2^{n^4/24-2n^3}$. If the sum of the absolute values of
all other primitive-pole coefficients is at most one quarter of
that amplitude, then a positive witness occurs in every window
of at most nine integers and both strict indices exceed
$2^{cn}$ for every fixed $c<1/4$. This total-remainder estimate
is not proved. Equivalently the signed wave is an explicit
coefficient of a finite negacyclic inverse, by (P-F56z);
its coefficients have no assumed sign. This route needs no
estimate for $B_n$ or for the other Hilbert poles.
The new comparison (P-F56ak)--(P-F56au) proves that every
normalized multiplier range $3\le J\le2^{\delta n}$ is
negligible for each fixed $\delta<1/600$, with an explicit
bound for $\delta=1/100000$. Only the multipliers outside
such a range remain in this sufficient estimate.
Alternatively (P-F56aa)--(P-F56aj) reduces the full mass
bound to a uniform two-child contraction for increment
products having $\binom{n-1}{2}$ factors. The error from
the two selected sibling branches is already bounded by
$2^{-n^3/12}$ for $n\ge48$. The uniform contraction itself
is unproved; a largest-child estimate alone would not suffice.
P-F47 identifies exterior-volume powers with nonvanishing at the
Fermat cubic. P-F49 constructs such invariants in polynomial
degree, with a sharper explicit bound from the supplied Dirichlet
theorem. P-F50 gives a polynomial eventual-support threshold on
every stretch allowed by that point's period. Combining this
with the odd top-exterior-power invariant gives polynomial
full-ring eventual support on the infinite family in
(P-F50s)--(P-F50t). These results
control realized point evaluations; the source's strict indices
also test constituent polynomials away from their residue classes.
P-F54 now proves unconditional polynomial eventual actual support
for every $n\ge33$, by a specified exactly balanced cubic and
quantified compact averaging. Every coefficient degree
$d\ge(10n)^{200}$ with $n\mid3d$ occurs at this point. Thus
the missing strict-index estimate cannot be attributed to an
unproved eventual-support bound on this cubic family; its
remaining content is the off-residue constituent tests.
Finite-dimension absorption in (P-F54y) gives an absolute,
not numerically specified, constant $C$ with full support for
all $n\ge4$ and $s\ge Cn^{199}$. The quotient surfaces in
(P-F54v)--(P-F54x) have every actual degree, fixed dimension
and $a$-invariant, and arbitrarily large strict indices.
P-F48 proves that the cubic family and its natural determinant
twists are exponentially sparse in each of the two stated SH(b)
input populations, so this amplification does not settle SH(b).
P-F38 excludes deriving this bound from small matrix entries, bounded
coordinates and cubic-scale dimension alone, even after normalizing by
the row rank. The actual cubic weight system and Weyl alternation remain
necessary extra structure. P-F33 shows why normal rational Gorenstein geometry with fixed
negative $a$-invariant and fixed dimension cannot alone supply
such a bound: its genuine arbitrarily large periods coexist with
zero indices and an unbounded average constant term. P-F35 gives
the exact calibration $p(f_4)=s(f_4)=23$ and minimal period $120$.

\item[O-COUNT: uniform effective shifted counting.]
P-F2 constructs a single shifted polytope from explicit positive Hilbert
data. P-F12 controls their existence and size under a bounded
positive-form premise in expanded parameters. P-F31 shows that the
proposed numerical PH3 bound cannot supply that premise in general. P-F13 reconstructs
and verifies candidate data from a finite prefix under an explicit
denominator promise. A procedure constructing them with the original
partition-input complexity bounds remains to be given. A small object
whose existence is known separately for every instance is not automatically
a polynomial-time constructible object. The homogeneous version is already
refuted by P-K1; no revision should silently restore that false assertion.

\item[O-NSQ: nonstandard geometric realization.]
For the particular bases and products in GCT VII--VIII, one must construct
the category, identify the basis, control parity or shifts, establish
product and specialization compatibility, and supply size bounds.
P-F6 is a sufficient criterion after these data exist. The determinant
and theta-lift passages of the supplied note do not themselves identify
the relevant Hom spaces. P-F5 refutes the simplest attempted substitution
of ordinary interlacing for plethysm restriction.
P-F36 additionally proves that the actual cubic example has a
quadratic coefficient extension that survives every finite power
regrading. Its separate conjugate blocks cannot descend through
an ordinary rational projector, and its displayed raw coefficient
cannot be a finite graded Laurent character. A comparison must
retain that extension and its descent datum, or explicitly recover
the blocks from a conjugate sum. Arbitrary coefficient systems
and the source's relaxed minimal-polynomial conjectures remain
outside this obstruction.
P-F39 supplies an actual rational symmetrizable positive operator
for the two conjugate scalars, the exact source normalization of
its trace and norm, and positive Tate--Lefschetz models for those
two normalized coefficients. Before normalization its Laurent
coefficients have negative entries. P-F52 now constructs the
missing actual-algebra intertwiner for this operator: the short
central word $(Q_1-Q_2)^2+S-f^2$ gives precisely its matrix,
and an explicit positive tensor form proves positivity at every
real $q>0$. The complete relation and cell-element factorization
prove the scalar equation independently of the source's
inconsistent printed two-by-two generator matrices. Comparison
of the coefficient models with the geometric or theta objects,
their products, and the prescribed normalization remains open.
The explicit alternative normalization in (P-F52m)--(P-F52p)
shows that positive actual action and positive Lefschetz trace
and norm do not even force the source factor to divide the scale.
The source's actual spectral-interpolation rule does give an
exact positive factorization of that scale in (P-F52q)--(P-F52s).
Moreover (P-F52t)--(P-F52x) refute two literal common-object
dictionaries: the trace and norm cannot be the first two ordinary
lambda characters under any common scalar normalization, and
the prescribed normalized pair cannot be simultaneous traces of
one Frobenius and its exterior square over every finite extension.
These are precise excluded comparisons, not a general exclusion
of geometric or theta constructions.
The positive order in (P-F52y)--(P-F52ah) now supplies a finite
cohomological additive endofunctor whose Grothendieck-group action
is explicitly intertwined with the actual central operator.
Its basis is an integral bar-fixed lattice in the source's
convenient word form. The corresponding global generator fails
membership in two explicitly rescaled generator forms, and its
change to the displayed source candidate basis has a negative
coefficient. Thus full source-canonical products, the ambient
minimum-degree condition, monoidal structure and uniform tensor
degree are the remaining comparisons; the finite central
intertwiner and additive coefficient model are already constructed.
The extension (P-F52av)--(P-F52bc) gives an actual
four-dimensional source left ideal with an integral, bar-fixed
lattice on which both chosen rescaled generators $FQ_1,FQ_2$
and all their positive words have nonnegative Laurent matrices.
The four explicitly specified source vectors give the full
intertwiner, and finite additive functors realize these word
actions on Grothendieck groups. Coherent source relations,
canonical products and the other remaining structures are
still additional requirements. Both complementary pairs
$FQ_i,FP_i$ cannot be simultaneously nonnegative in one
basis where the actual generators do not commute.
The quadratic functor relations now have explicit natural
isomorphisms coherent on repeated equal generators.
The ordinary direct-sum model does not make these functors
self-adjoint, but (P-F52bg)--(P-F52bl) supplies a nonnegative
Laurent Gram matrix intertwining both adjoints and positive
definite at every positive parameter other than one. Thus
numerical adjoint compatibility is available; actual Hom
composition, adjunction and mixed-relation coherence remain open.

\item[O-THETA: comparison with an identified invariant.]
P-F37 supplies the actual meromorphic automorphic ring of
five-variable cubics. P-F41 identifies its weight-$48$
Borcherds form with the ordinary degree-$80$ cubic discriminant,
including all products in the generated polynomial subalgebra.
The specified construction uses a regularized multiplicative
orthogonal theta lift, an $E_6$ quasi-pullback and restriction.
The next comparison must identify this named output with the
supplied iterated unitary lift or determinant, or produce an
identified larger graded collection with quantitative bounds.
The supplied note already mentions orthogonal groups; the group
name alone neither establishes nor excludes this comparison.
P-F46 identifies the exact positive input as the vector-valued
series $\Theta_{E_6}/\delta$, including the discriminant-group
gluing and principal coefficients $1,96,27,27$. Its ordinary
Petersson norm diverges. Its positive Fourier coefficient counts
grow faster than the actual $T^{(5)}$ Hilbert coefficients under
every positive affine reindexing. Thus the remaining comparison
must identify the regularized multiplicative operation or another
explicit coefficient-space construction; direct reindexing of
these input counts is refuted. Hofmann's unitary Borcherds
construction identifies the same form as a multiplicative
unitary lift of this input; the remaining comparison concerns
that exact regularized lift and the supplied iterated definite
unitary operations, not whether the output can be called unitary.
P-F51 further realizes the exact positive $E_6$ factor as an
ordinary convergent theta integral from a specified nonnegative
function on a definite $U(3)$ quotient. Its lattice cosets,
finite level, splitting on $U(3)\times SL_2$, individual-class
measure and $y^{-3/2}$ normalization are fixed. The compatible
right unitary group is split $U(1,1)$. The remaining comparison
must match the supplied definite chains to these specified
vectors and account for division by $\delta$ and the regularized
multiplicative operation; the positive factor itself is identified.
P-F53 realizes all coefficients of that exact input by finite
disjoint unions of K3 Hilbert schemes, with even cohomology
and purity in each cohomological degree. Its completed
lambda-ring identity supplies a geometric source for the
inverse-discriminant coefficient series. Its extension
(P-F53h)--(P-F53p) also identifies the known normalized cusp
product with the completed exterior Euler character of those
spaces and recovers every input coefficient by a formal inverse.
No finite actual or virtual positive-height object gives that
same exterior character. The comparison with the supplied
iterated operation and with all actual multiplicity maps is
still missing; the completed character-level operation is fixed.
P-F55 identifies the square of the same orthogonal theta output
with the six-variable cubic discriminant, of coefficient degree
$192$, and matches its cyclic restriction with the square of
the five-variable discriminant. This adds an exact fixed-dimensional
compatibility; it does not supply the missing comparison with
the supplied iterated operations or an asymptotic family.
P-F45 proves that the full cyclic-cubic-cover eigenspace
classifying domain is a ball exactly for $n=4,5$. Beyond
this range an asymptotic construction must specify what
replaces that full-domain model; selected substructures,
special period images and other theta constructions remain open.

\item[O-OBSTRUCTION: an actual family of separating equations.]
In the no-occurrence range, prove a strict inequality between the actual
coordinate-ring multiplicities, or construct an explicit polynomial
vanishing on the determinant closure and not on the padded permanent.
The ambient plethysm and Kronecker multiplicities do not determine the
ranks of the quotient maps. P-F8 and PL-06 isolate this rank/comparison
requirement. No such asymptotic obstruction family is constructed here.

\item[O-DIR-HN: arbitrary solution fields in Nullstellensatz.]
C-GRH-01 proves only the cyclotomic-field transfer. Koiran's algorithm uses
a uniform prime-ideal/root-modulo-prime estimate for arbitrary primitive
solution fields. Rojas's DZH or GIPIT conditions are alternative sufficient
analytic inputs. The present supplied Dirichlet theorem does not state
these inputs. No formal nonimplication between the analytic conjectures
is proved.

\item[O-PIT: the derandomization required by GCT V.]
The exact algorithm or hardness hypothesis of each theorem must still be
provided. The supplied counting algorithm yields C-PIT-01 in a bounded
finite-field circuit model. It does not supply the characteristic-zero
black-box, border or strengthened symbolic determinant PIT assertions
used elsewhere. GRH does not replace those separate hypotheses.

\item[O-INTEGRAL: uniform bounds for specialization.]
B7.1 proves good-prime survival for a fixed integral separator; B7.2 gives
an exact rank test in any fixed degree and integer minors excluding bad
primes. An efficient GCT transfer needs a uniform separator-degree bound,
succinct rank data and a usable prime bound. Neither Hilbert-scheme
representability nor the pointwise specialization theorem supplies these
quantitative statements.

\item[O-COMPLEXITY: Boolean and algebraic separation.]
The ETH theorem refutes the precise exponential hypotheses treated in
C-SAT-01. Its supplied circuit corollary is superlinear for
$\mathsf E^{\mathsf{NP}}$, not superpolynomial for $\mathsf{NP}$.
The current constant-free integer Boolean--arithmetic disjunction cited
in the complexity section does not choose a branch and does not identify
the usual complex-field $\mathsf{VP}$ with its constant-free integer
variant. The stated upper bounds and conditional bridges do not settle
$\mathsf P$ versus $\mathsf{NP}$ or the target GCT obstruction.
\end{description}

\section{Approaches tried and the next discriminating steps}
The qualitative-sign route was tested by actual determinant projectors
and a virtual-character counterexample (P-F3--P-F4). Its next step is an
identity with the intended multiplicity projector together with a size
bound. The standard-branching route was tested on $\operatorname{Sym}^2$
(P-F5); a replacement must use the actual Schur-functor homomorphism.

The purity route was tested at three increasingly specific levels:
semisimple categorical structure (P-F6), pure finite schemes encoding
\#3SAT (P-F7), and the actual canonical plethysm algebra carrying the
obstructing Hilbert function, a positive multiplicative basis and pure
intersection cohomology (PLE-F7--PLE-F8). The next step is an
equivariant graded comparison theorem, not another proof that an unrelated
cohomology group is pure. P-F10 separately tests the missing quantitative
control of a grading.

The Hilbert-series route made positive progress: P-F1 constructs the
optimal shifted model on the obstructing ray; PLE-F5 realizes its normal
invariant ring; PLE-F6 classifies all rescalings that repair homogeneous
Ehrhart realization; P-F2 treats every explicit positive rational form;
P-F9 and P-F11 supply coefficient-positivity bounds; P-F12 connects the
last bound to the numerical positive-form clause of revised PH3.
P-F13 gives finite reconstruction under a denominator promise;
P-F16--P-F18 give finite fixed-embedding and fixed-dimension bounds;
P-F19 refutes polynomial growth of the canonical parameter bound.
P-F30--P-F31 subsequently close the numerical positive-form question
negatively, using a complete support theorem and a fixed-component
calculation. The next step is a direct PH2 or saturation argument,
together with uniform construction of a suitable revised model.
The ordinary pole is proved by these additional arguments, not
inferred from the h.s.o.p. bound alone.
P-F20 confirms that distinction with universal denominators and
fixed-target covariants. Its multivariate counterexample defeats a
transfer based only on quasi-polynomiality from an affine slice to
simultaneous rays. P-F21--P-F22 then defeat generic normal,
Cohen--Macaulay, rational and Gorenstein attempts to repair the
parameter-degree inference. P-F23--P-F24 made the geometric gate precise: the known chain
component is smooth, and the other fixed components had to be
controlled. P-F26 identifies the entire smooth fixed locus with
the chain orbit; P-F27 excludes every graph-support boundary
competitor, including rank-deficient supports and lifts of larger
two-power order. P-F29 handles the complementary cube-and-squarefree
case. P-F30 excludes all remaining mixed supports, and P-F31 applies
the noncancellation criterion. No competing nonzero component remains. The binary determinant subfamily is separately treated
in P-F25: its realized holes are bounded by thirteen, but the
source's strict saturation definition also tests off-residue
constituent values and is not settled by that table.

The archive route distinguishes the three products on representations,
separates the two Burnside operations, checks the transcribed identities
on small groups, and treats the modular and quadratic-form caveats
(PL-01--PL-08). Its obstruction gate is the actual rank of the
Foulkes--Howe multiplicity-space map. A Schur-positive character of its
source and target does not compute that rank.

The norm and approximation routes were tested against scaling, exact
zero detection, adjacent integer multiplicities, and an explicit rank-two
family approaching rank one in normalized norm (B4.1--B5.3). They do give
valid contraction-model obstructions under the stated operator bounds.
A full algebraic-orbit obstruction requires an invariant and a certified
separating gap that survive the allowed degenerations.

The arithmetic and ETH routes track the exact analytic family, all input
and reduction sizes, the direction of reductions, arithmetic precision,
and the class in which a lower bound is obtained. Their next steps are
specified in O-DIR-HN, O-PIT and O-COMPLEXITY. These are not resolved by
relabeling the supplied subexponential bound as polynomial time.

The exterior-volume route is now an exact Fermat-evaluation
dictionary (P-F47). A prime-characteristic determinant identity,
applied through normalized Hessians and Cayley's contraction,
gives the actual invariants of P-F49--P-F50. The remaining
strict-index step cannot be replaced by this point-specific
conductor. P-F48 separately closes the proposed density
amplification by ordinary determinant twists, with the exact
rank and partition-size conventions of each source clause.
